\documentclass[11pt, a4paper]{article}

% --- OVERLEAF COMPATIBLE PREAMBLE ---
\usepackage[utf8]{inputenc}
\usepackage[T1]{fontenc}
\usepackage{lmodern}
\usepackage[a4paper, top=2.5cm, bottom=2.5cm, left=2cm, right=2cm]{geometry}
\usepackage[english]{babel}
\usepackage{amsmath, amssymb}
\usepackage{booktabs}
\usepackage{enumitem}
\usepackage{microtype}
\usepackage[colorlinks=true, linkcolor=blue, citecolor=blue, urlcolor=blue]{hyperref}

\setlist[itemize]{leftmargin=*, noitemsep, topsep=3pt}
\setlist[enumerate]{leftmargin=*, noitemsep, topsep=3pt}

\title{\textbf{Graph-Augmented Entity Resolution and Topological Reasoning in Historical Genealogical Networks}}
\author{\textbf{Rory Macdonald} \\ Course Project Proposal (Milestone 1) \\ CISC 451/839: Topics in Data Analytics}
\date{\today}

\begin{document}

\maketitle

\noindent\textbf{Project Category:} Category 1 (Improving Existing Approaches) / Category 2 (Exploratory)

\vspace{0.8em}

\section{Problem Statement}
Historical genealogical databases compiled across archival collections frequently contain fragmented, duplicate, or inconsistent entries. Standard record linkage techniques rely heavily on localized string similarity and attribute pooling (e.g., grouping by exact names). These methods fail when naming conventions vary drastically (e.g., Acadian \textit{dit} names) or when critical attributes like birth dates or locations are missing or misspelled. 

This project investigates how relational structural context (family tree graph neighborhoods) and Graph-Augmented Retrieval and Reasoning can be combined to resolve entity duplication across multi-tree genealogical networks.

\begin{itemize}
    \item \textbf{Research Question 1 (RQ1):} How effectively do topological Graph Neural Network (GNN) neighborhood embeddings improve duplicate entity candidate retrieval compared to text/attribute-only vector embeddings on sparse historical genealogical trees?
    \item \textbf{Research Question 2 (RQ2):} How can Graph Retrieval-Augmented Generation (GraphRAG) and step-by-step reasoning traces improve the precision and interpretability of candidate duplicate verification for domain practitioners?
\end{itemize}

\section{Challenges}
\begin{enumerate}
    \item \textbf{Graph Sparsity \& Structural Heterogeneity:} Historical family trees vary widely in depth, missing nodes, and incomplete branch linkages, making uniform topological graph embedding learning challenging.
    \item \textbf{Scalability of Subgraph Comparison:} Performing pairwise topological neighborhood matching across a large database (150,000+ nodes) requires efficient multi-hop neighborhood extraction and vector indexing mechanisms.
    \item \textbf{Lack of Ground-Truth Benchmarks:} Historical archival datasets lack exhaustively labeled duplicate pairs, requiring robust synthetic noise-injection frameworks for systematic evaluation.
\end{enumerate}

\section{Dataset}
\begin{itemize}
    \item \textbf{Acadian Genealogy Database:} A historical dataset comprising over 150,000 individual records across 60 interconnected family trees. Each record contains attributes (names, dates, and locations for birth, baptism, marriage, and death) alongside relational links (parents, children, and spouses).
    \item \textbf{Evaluation Ground Truth:} Generated by extracting verified multi-tree sub-lineages and introducing synthetic perturbations (e.g., spelling variations, omitted dates, missing parent links) alongside expert-verified duplicate records to benchmark candidate retrieval.
\end{itemize}

\section{Literature Review}
The project is informed by three key papers spanning graph representation learning, advanced retrieval, and graph-augmented reasoning:
\begin{enumerate}
    \item \textbf{Edge et al. (2024) -- \textit{From Local to Global: A Graph RAG Approach to Query-Focused Summarization}:} Demonstrates how extracting local graph structures and community hierarchies improves context retrieval over standard vector search.
    \item \textbf{Tian et al. (2024) -- \textit{Graph Neural Prompting with Large Language Models}:} Establishes frameworks for translating GNN-encoded structural graph embeddings into promptable context for language models to perform downstream relational reasoning.
    \item \textbf{Marisca et al. (2022) -- \textit{Learning to Reconstruct Missing Data from Spatiotemporal Graphs with Sparse Observations} (NeurIPS):} Provides foundational approaches for processing incomplete, sparse structural networks with graph neural architectures.
\end{enumerate}

\section{Planned Approach}
We propose a two-stage pipeline combining graph machine learning with agentic reasoning:
\begin{itemize}
    \item \textbf{Stage 1: Topological Candidate Retrieval (RQ1):} We construct a relational kinship graph where nodes represent individuals and edges represent relational links (parent-child, spouse). We train a Relational Graph Convolutional Network (RGCN) to output dense node embeddings encoding both attribute data and multi-hop neighborhood topology. We evaluate candidate retrieval against an attribute-only baseline (concatenated text fields passed through a text embedder) using vector cosine similarity.
    \item \textbf{Stage 2: GraphRAG Verification \& In-Text Reasoning (RQ2):} For top-$k$ retrieved candidates, we extract local 2-hop subgraphs (parents, siblings, spouses, spatiotemporal markers). We feed these subgraphs into a Language Model to generate explicit, step-by-step reasoning traces evaluating chronological consistency, kinship overlap, and identity alignment to produce a final classification.
\end{itemize}

\section{Evaluation}
\begin{itemize}
    \item \textbf{Quantitative Retrieval Performance (RQ1):} Evaluated using Recall@$K$ (is the true duplicate in the top 5 or 10 retrieved candidates?), Mean Reciprocal Rank (MRR), and Precision-Recall Curves compared directly against attribute-only vector baselines and exact-name pooling.
    \item \textbf{Verification \& Reasoning Quality (RQ2):} Measured by overall Precision, Recall, and F1-score on duplicate classification before and after the GraphRAG verification step. Faithfulness and logical consistency of generated reasoning traces will be assessed using RAGAS (Retrieval-Augmented Generation Assessment) context precision metrics to ensure explanations are strictly grounded in graph context without hallucination.
\end{itemize}

\end{document}